\section{Introduction}
\label{sec:intro}

Compiler pipelines are increasingly built and operated by automated agents, yet the
frameworks they must extend were designed for humans editing C++ and waiting for
rebuilds. Adding a pass, an operation definition, a verifier, or an invariant to an
\mlir{}-based stack today means editing TableGen and C++ sources, rebuilding the tool,
and restarting the process. That authoring loop --- not the IR data structure --- is
the bottleneck.

Our thesis is that the IR \emph{core} is not the right place to compete: \mlir{} has
won the core, and systems that try to replace it lose the ecosystem while chasing a
moving target. What is missing around the core are (i) declarative extension tiers
that an agent can author as data, (ii) a control plane that discovers, validates, and
loads those extensions at runtime, and (iii) first-class invariants that make
agent-authored transformations trustworthy.

\textbf{Contributions (draft).}
\begin{itemize}
  \item A tiered extension model for agent-authored compiler extensions --- from
        pipeline manifests and interpreted transformations up to runtime operation
        definitions and compiled plugins --- with precise rebuild boundaries.
        {\color{red}\textbf{TODO}}
  \item A runtime control plane: a manifest registry, plugin loader, and structured
        (machine-readable) diagnostics.
        {\color{red}\textbf{TODO}}
  \item Invariants as data: per-stage contracts checked at runtime.
        {\color{red}\textbf{TODO:} decide SMT-checked vs.\ translation-validation depth.}
  \item An evaluation protocol for \emph{agent operability}: task success, time and
        cost to first successful extension, error recovery, versus a vanilla
        out-of-tree \mlir{} workflow.
        {\color{red}\textbf{TODO:} define tasks, baselines, metrics.}
\end{itemize}
